% Blank assessment / workbook variant of src/security/application-security/processes/application-review-analysis-template.tex
% No completed results or approval implied. See authoring source for guidance.
\documentclass[10pt]{article}
\usepackage[letterpaper,margin=0.72in,headheight=15pt]{geometry}
\usepackage[T1]{fontenc}
\usepackage[utf8]{inputenc}
\usepackage{lmodern}
\usepackage{microtype}
\usepackage{parskip}
\usepackage{enumitem}
\usepackage{booktabs}
\usepackage{tabularx}
\usepackage{longtable}
\usepackage{array}
\usepackage{xcolor}
\usepackage[most]{tcolorbox}
\usepackage{fancyhdr}
\usepackage{hyperref}
\usepackage{amsmath,amssymb}
\usepackage{listings}
\usepackage{tikz}
\usetikzlibrary{arrows.meta,positioning,shapes.geometric,fit}
\usepackage{needspace}
\usepackage{multicol}
\usepackage{lastpage}
\usepackage{titlesec}

\definecolor{navy}{HTML}{17324D}
\definecolor{bluegray}{HTML}{43657D}
\definecolor{lightblue}{HTML}{EAF2F8}
\definecolor{lightgray}{HTML}{F3F5F7}
\definecolor{linegray}{HTML}{C9D1D9}
\definecolor{darkgray}{HTML}{4A5561}
\definecolor{accent}{HTML}{8E3B46}

\hypersetup{colorlinks=true,linkcolor=navy,urlcolor=bluegray,citecolor=navy,pdfborder={0 0 0}}
\setlist[itemize]{leftmargin=1.4em,itemsep=2pt,topsep=4pt}
\setlist[enumerate]{leftmargin=1.6em,itemsep=2pt,topsep=4pt}
\renewcommand{\arraystretch}{1.22}
\setlength{\tabcolsep}{4pt}
\newcolumntype{Y}{>{\raggedright\arraybackslash}X}
\newcolumntype{C}[1]{>{\centering\arraybackslash}p{#1}}
\newcolumntype{L}[1]{>{\raggedright\arraybackslash}p{#1}}
\pagestyle{fancy}
\fancyhf{}
\lhead{Application Security Engineering}
\rhead{Code Review Handbook}
\cfoot{\thepage\ of \pageref{LastPage}}
\renewcommand{\headrulewidth}{0.4pt}
\titleformat{\section}{\Large\bfseries\color{navy}}{\thesection}{0.65em}{}
\titleformat{\subsection}{\large\bfseries\color{bluegray}}{\thesubsection}{0.55em}{}
\titleformat{\subsubsection}{\normalsize\bfseries\color{darkgray}}{\thesubsubsection}{0.5em}{}
\titlespacing*{\section}{0pt}{1.4em}{0.55em}
\titlespacing*{\subsection}{0pt}{1.0em}{0.35em}

\newcommand{\checkbox}{\ensuremath{\square}}
\newcommand{\fieldline}[1]{\textbf{#1}\hspace{0.5em}\leaders\hbox{\rule[0.4ex]{0.6em}{0.35pt}\hspace{0.25em}}\hfill\kern0pt}
\newcommand{\workarea}[1]{\par\smallskip\noindent\begin{tcolorbox}[height=#1,colback=white,colframe=linegray,boxrule=0.5pt,arc=0pt,left=2pt,right=2pt,top=2pt,bottom=2pt]\mbox{}\end{tcolorbox}\par}
\newtcolorbox{keybox}[1]{colback=lightblue,colframe=bluegray,title=\textbf{#1},fonttitle=\bfseries,boxrule=0.7pt,arc=2pt,left=6pt,right=6pt,top=5pt,bottom=5pt}
\newtcolorbox{cautionbox}[1]{colback=lightgray,colframe=accent,title=\textbf{#1},fonttitle=\bfseries,boxrule=0.7pt,arc=2pt,left=6pt,right=6pt,top=5pt,bottom=5pt}
\newtcolorbox{worksheetbox}[1]{colback=white,colframe=linegray,title=\textbf{#1},fonttitle=\bfseries,coltitle=navy,boxrule=0.6pt,arc=1pt,left=6pt,right=6pt,top=5pt,bottom=5pt}
\lstset{basicstyle=\ttfamily\small,columns=fullflexible,breaklines=true,frame=single,framerule=0.4pt,rulecolor=\color{linegray},backgroundcolor=\color{lightgray},showstringspaces=false,tabsize=2,keywordstyle=\bfseries}
\setcounter{tocdepth}{2}

\title{\textbf{Application Review Analysis --- Blank}\\\large Unfilled Application Security Assessment Workbook}
\author{}
\date{}

\begin{document}
\maketitle
\thispagestyle{empty}
\begin{keybox}{Purpose}
Blank engagement worksheet for authorized application security review. Complete each field with project-specific evidence; no review results are implied.
\end{keybox}
\vfill
\noindent\textbf{Use note.} This handout is a practitioner aid for authorized software security assessment. It emphasizes evidence, traceability, and disciplined review rather than a single rigid methodology.
\newpage
\tableofcontents
\newpage

\section{Operating Model}
\noindent Record scope and authorization before assessment. Complete each
subsequent worksheet with directly observed evidence, rather than assumed findings.

\section{Document Control}
\begin{worksheetbox}{Engagement Record}
\par\smallskip\noindent\begin{tabularx}{\linewidth}{@{}L{0.26\linewidth}Y@{}}
Application / system name & \\ \midrule
Application owner & \\ \midrule
Business owner & \\ \midrule
Development team & \\ \midrule
Assessment lead & \\ \midrule
Additional reviewers & \\ \midrule
Assessment dates & \\ \midrule
Review type & \\ \midrule
Repository / repositories & \\ \midrule
Version / commit / build & \\ \midrule
Deployment environment & \\ \midrule
Assessment status & \\ \midrule
Classification & \\ \midrule
Review constraints & \\ 
\end{tabularx}
\end{worksheetbox}

\section{Review Objectives and Scope}
\begin{worksheetbox}{Objective Statement}
\par\smallskip\noindent\begin{tabularx}{\linewidth}{@{}L{0.29\linewidth}Y@{}}
Assessment purpose & \\ \midrule
Business drivers & \\ \midrule
Security objectives & \\ \midrule
Success criteria & \\ \midrule
Expected depth & \\ \midrule
Expected coverage & \\ \midrule
Known constraints & \\ \midrule
Explicit exclusions & \\ \midrule
Regulatory / policy drivers & \\ \midrule
Primary stakeholders & \\ \midrule
Audience for findings & \\ 
\end{tabularx}
\end{worksheetbox}

\subsection{Review Intent}
Select all that apply and define what ``enough'' means for this engagement.
\par\smallskip\noindent\begin{tabularx}{\linewidth}{@{}C{0.07\linewidth}L{0.34\linewidth}Y@{}}
\toprule
Select & Intent & Engagement-specific interpretation \\
\midrule
\checkbox & Rapid discovery of significant vulnerabilities & \\ \addlinespace[1.2em]
\checkbox & Broad security coverage & \\ \addlinespace[1.2em]
\checkbox & Detailed assurance of selected components & \\ \addlinespace[1.2em]
\checkbox & Targeted review of high-risk paths & \\ \addlinespace[1.2em]
\checkbox & Verification of prior remediation & \\ \addlinespace[1.2em]
\bottomrule
\end{tabularx}
\clearpage

\section{Application Access Profile}
Access determines which forms of evidence are available and places practical limits on achievable assurance. Treat these categories as planning aids, not immutable labels.
\begin{longtable}{@{}L{0.16\linewidth}L{0.20\linewidth}L{0.19\linewidth}L{0.20\linewidth}L{0.18\linewidth}@{}}
\toprule
Mode & Available evidence & Typical analysis & Important limitation & Compensating action \\
\midrule
\endfirsthead
\toprule
Mode & Available evidence & Typical analysis & Important limitation & Compensating action \\
\midrule
\endhead
Source only & Source tree and documentation & Static review, search, source navigation, manual flow analysis & Runtime behavior and environment assumptions may be unverified & Reconstruct build/runtime assumptions; flag dynamic questions explicitly \\
\addlinespace
Binary only & Executable artifacts & Dynamic analysis, debugging, reverse engineering & Intent and high-level structure are harder to infer & Build behavioral model; preserve runtime traces and call evidence \\
\addlinespace
Source + runtime & Source plus functional build / test system & Combined static and dynamic review & Test environment may differ from production & Record environment differences and validate critical assumptions \\
\addlinespace
Checked / debug build & Binary with extra symbols or diagnostics & Dynamic tracing with richer context & Still lacks full source intent & Map symbols to behavior; retain uncertainty about source-level semantics \\
\addlinespace
Strict black box & External interfaces only & Dynamic testing, fuzzing, protocol / API observation & Limited implementation visibility places an upper bound on coverage claims & State the assurance limit; use reproducible external evidence \\
\bottomrule
\end{longtable}

\begin{worksheetbox}{Selected Access Profile}
\fieldline{Primary mode:}\par
\fieldline{Available artifacts:}\par
\fieldline{Unavailable artifacts:}\par
\fieldline{Analysis methods available:}\par
\fieldline{Assurance implications:}\par
\fieldline{Compensating techniques:}
\end{worksheetbox}
\clearpage

\section{Timeline and Coverage Planning}
Observed review velocity varies with language, layering, code quality, familiarity, and the review objective. Track actual progress rather than assuming a universal lines-per-hour rate.

\begin{worksheetbox}{Planning Baseline}
\par\smallskip\noindent\begin{tabularx}{\linewidth}{@{}L{0.32\linewidth}Y@{}}
Total review time & \\ \midrule
Reviewer count & \\ \midrule
Estimated codebase size & \\ \midrule
Languages / frameworks & \\ \midrule
Application complexity & \\ \midrule
Documentation quality & \\ \midrule
Runtime / QA availability & \\ \midrule
Desired coverage & \\ \midrule
High-risk areas & \\ \midrule
Excluded areas & \\ \midrule
Review milestones & \\ 
\end{tabularx}
\end{worksheetbox}

\subsection{Observed Review Velocity}
\par\smallskip\noindent\begin{tabularx}{\linewidth}{@{}L{0.18\linewidth}L{0.18\linewidth}L{0.17\linewidth}L{0.17\linewidth}Y@{}}
\toprule
Date / iteration & Component & Time spent & Approx. scope & What affected pace? \\
\midrule
\rule{0pt}{2.4em} & & & & \\
\rule{0pt}{2.4em} & & & & \\
\rule{0pt}{2.4em} & & & & \\
\rule{0pt}{2.4em} & & & & \\
\bottomrule
\end{tabularx}
\clearpage

\section{Information Collection}
Build the first working model from multiple evidence sources. Treat documents and interviews as claims to validate against implementation when the review reaches the relevant code.

\begin{longtable}{@{}L{0.27\linewidth}C{0.09\linewidth}L{0.22\linewidth}L{0.15\linewidth}L{0.20\linewidth}@{}}
\toprule
Evidence source & Avail.? & Location / reference & Owner & Reliability / notes \\
\midrule
\endfirsthead
\toprule
Evidence source & Avail.? & Location / reference & Owner & Reliability / notes \\
\midrule
\endhead
Developer interviews & & & & \\
Architecture documentation & & & & \\
Design documentation & & & & \\
Standards / protocol specifications & & & & \\
API documentation & & & & \\
Threat models & & & & \\
Data-flow diagrams & & & & \\
Source profiling & & & & \\
Dependency inventory & & & & \\
System / deployment profiling & & & & \\
Build instructions & & & & \\
Test cases & & & & \\
Prior findings & & & & \\
Previous penetration tests & & & & \\
Operational controls & & & & \\
\bottomrule
\end{longtable}
\clearpage

\section{Application Model}
Start broad. Refine the model as code evidence accumulates. In a bottom-up or hybrid review, incomplete models are expected; record uncertainty instead of hiding it.

\subsection{General Purpose}
\rule{0.82\linewidth}{0.4pt}
\workarea{1.0in}

\subsection{Assets}
\rule{0.82\linewidth}{0.4pt}
\workarea{1.0in}

\subsection{Entry Points}
\par\smallskip\noindent\begin{tabularx}{\linewidth}{@{}L{0.20\linewidth}L{0.20\linewidth}L{0.20\linewidth}Y@{}}
\toprule
Entry point & Input / identity & Trust level & Security relevance / notes \\
\midrule
\rule{0pt}{2.6em} & & & \\
\rule{0pt}{2.6em} & & & \\
\rule{0pt}{2.6em} & & & \\
\rule{0pt}{2.6em} & & & \\
\bottomrule
\end{tabularx}

\subsection{Components, Relationships, and Trust Boundaries}
\par\smallskip\noindent\begin{tabularx}{\linewidth}{@{}L{0.22\linewidth}L{0.20\linewidth}L{0.22\linewidth}Y@{}}
\toprule
Component / module & Role & Interfaces / relationships & Trust / security expectation \\
\midrule
\rule{0pt}{2.8em} & & & \\
\rule{0pt}{2.8em} & & & \\
\rule{0pt}{2.8em} & & & \\
\bottomrule
\end{tabularx}

\subsection{External Dependencies}
\par\smallskip\noindent\begin{tabularx}{\linewidth}{@{}L{0.24\linewidth}L{0.18\linewidth}L{0.23\linewidth}Y@{}}
\toprule
Dependency & Type & Security assumption & Validation / evidence \\
\midrule
\rule{0pt}{2.7em} & & & \\
\rule{0pt}{2.7em} & & & \\
\rule{0pt}{2.7em} & & & \\
\bottomrule
\end{tabularx}

\section{Review Structure Selection}
\par\smallskip\noindent\begin{tabularx}{\linewidth}{@{}L{0.16\linewidth}L{0.26\linewidth}L{0.24\linewidth}Y@{}}
\toprule
Structure & Best fit & Main advantage & Main risk / switch signal \\
\midrule
Top-down & Reliable design and threat-model information exists & Directs effort toward known security-relevant paths and design concerns & Stale or incomplete design can hide relevant paths; switch when implementation diverges materially \\
\addlinespace
Bottom-up & Design information is weak or unavailable & Builds understanding from implementation evidence and can expose undocumented behavior & Slow; may spend time on low-value code before relevance is understood \\
\addlinespace
Hybrid & Evidence quality varies across the application & Moves between design and implementation as knowledge improves & Requires disciplined notes so assumptions and structure changes remain visible \\
\bottomrule
\end{tabularx}

\begin{worksheetbox}{Selection Record}
\fieldline{Selected review structure:}\par
\fieldline{Why this structure fits the current evidence:}\par
\fieldline{Conditions that would trigger a change:}
\end{worksheetbox}
\clearpage

\section{Iterative Review Loop}
Use review iterations short enough to preserve focus and long enough to complete a meaningful target. A target is usually a bounded component, path, algorithm, candidate set, or design question.

\subsection{PLAN Worksheet}
\par\smallskip\noindent\begin{tabularx}{\linewidth}{@{}L{0.29\linewidth}Y@{}}
Current objective & \\ \midrule
Selected audit strategy & \\ \midrule
Why this strategy now? & \\ \midrule
Target component / path & \\ \midrule
Abstraction level & \\ \midrule
Expected duration & \\ \midrule
Assumptions & \\ \midrule
Candidate vulnerability classes & \\ \midrule
Required tools / evidence & \\ \midrule
Coordination / dependencies & \\ \midrule
Exit criteria & \\ 
\end{tabularx}

\subsection{WORK Worksheet}
\begin{longtable}{@{}L{0.30\linewidth}L{0.64\linewidth}@{}}
\toprule
Record & Notes / evidence reference \\
\midrule
\endfirsthead
\toprule
Record & Notes / evidence reference \\
\midrule
\endhead
Files / modules / functions reviewed & \\
Paths analyzed & \\
Evidence gathered & \\
Hypotheses & \\
Candidate findings & \\
Unexpected behavior & \\
Unanswered questions & \\
Working-paper references & \\
Coverage evidence & \\
Deferred investigation & \\
\bottomrule
\end{longtable}

\begin{cautionbox}{Time Management Guardrail}
If an interesting technical question begins consuming the iteration without improving coverage or a stated objective, record it in the hypothesis register and defer it deliberately. Deep dives are appropriate when the expected security value justifies the opportunity cost.
\end{cautionbox}
\clearpage

\subsection{REFLECT Worksheet}
\begin{itemize}
\item What did we learn about the application or its developers' implementation patterns?
\item Are we still examining the most security-relevant components?
\item Did any tangent consume more time than its expected value justified?
\item Did the application model change? Which evidence caused the change?
\item Are working papers sufficient for another reviewer to reconstruct the analysis?
\item Which assumptions were disproved, weakened, or strengthened?
\item Which hypotheses remain plausible?
\item Does the current strategy still fit the review stage and evidence?
\item Should the overall review structure change?
\item What should the next iteration target?
\end{itemize}
\fieldline{Next review action:}
\workarea{0.75in}

\section{Master Ideas / Hypothesis Register}
Capture speculative attack ideas quickly, then validate them when the review plan makes time for them.
\begin{longtable}{@{}C{0.06\linewidth}L{0.21\linewidth}L{0.14\linewidth}L{0.13\linewidth}L{0.17\linewidth}L{0.21\linewidth}@{}}
\toprule
ID & Hypothesis / attack idea & Component & Trust boundary & Validation approach & Status / finding \\
\midrule
\endfirsthead
\toprule
ID & Hypothesis / attack idea & Component & Trust boundary & Validation approach & Status / finding \\
\midrule
\endhead
\rule{0pt}{2.6em} & & & & & \\
\rule{0pt}{2.6em} & & & & & \\
\rule{0pt}{2.6em} & & & & & \\
\rule{0pt}{2.6em} & & & & & \\
\rule{0pt}{2.6em} & & & & & \\
\bottomrule
\end{longtable}
\clearpage

\section{Coverage Register}
Coverage evidence supports statements about what was reviewed; it does not prove the absence of vulnerabilities.
\begin{longtable}{@{}L{0.22\linewidth}L{0.18\linewidth}L{0.20\linewidth}L{0.11\linewidth}L{0.10\linewidth}L{0.11\linewidth}@{}}
\toprule
Component / path & Security relevance & Strategy / technique & Depth & Reviewer & Status / notes \\
\midrule
\endfirsthead
\toprule
Component / path & Security relevance & Strategy / technique & Depth & Reviewer & Status / notes \\
\midrule
\endhead
\rule{0pt}{2.5em} & & & & & \\
\rule{0pt}{2.5em} & & & & & \\
\rule{0pt}{2.5em} & & & & & \\
\rule{0pt}{2.5em} & & & & & \\
\rule{0pt}{2.5em} & & & & & \\
\bottomrule
\end{longtable}

\section{Working Papers Standard}
Working papers should be sufficiently clear that a qualified peer can approximate what was done, understand why a path was accepted or discarded, and trace a reported finding back to evidence.
\begin{multicols}{2}
\begin{itemize}
\item[\checkbox] review notes
\item[\checkbox] code locations
\item[\checkbox] hypotheses
\item[\checkbox] diagrams / models
\item[\checkbox] commands / searches
\item[\checkbox] tool output
\item[\checkbox] test inputs
\item[\checkbox] debugger observations
\item[\checkbox] candidate points
\item[\checkbox] discarded candidates and rationale
\item[\checkbox] coverage evidence
\item[\checkbox] peer-review notes
\end{itemize}
\end{multicols}
\clearpage

\section{Finding Analysis Template}
\begin{worksheetbox}{Finding Identity}
\par\smallskip\noindent\begin{tabularx}{\linewidth}{@{}L{0.24\linewidth}Y@{}}
Finding ID / title & \\ \midrule
Affected component & \\ \midrule
Repository / file & \\ \midrule
Function / method / lines & \\ \midrule
Vulnerability class & \\ 
\end{tabularx}
\end{worksheetbox}

\begin{worksheetbox}{Technical Analysis}
\fieldline{Technical description:}\workarea{0.8in}
\fieldline{Triggering conditions:}\workarea{0.55in}
\fieldline{Attacker-controlled inputs:}\workarea{0.45in}
\fieldline{Trust boundary / privilege transition:}\workarea{0.45in}
\fieldline{Prerequisites:}\workarea{0.45in}
\fieldline{Exploit / failure path:}\workarea{0.7in}
\fieldline{Security impact:}\workarea{0.55in}
\fieldline{Business impact:}\workarea{0.55in}
\end{worksheetbox}

\begin{worksheetbox}{Evidence and Resolution}
\fieldline{Evidence / working-paper references:}\workarea{0.55in}
\fieldline{Reproduction / validation:}\workarea{0.55in}
\fieldline{Compensating controls:}\workarea{0.45in}
\fieldline{Primary remediation:}\workarea{0.55in}
\fieldline{Remediation alternatives:}\workarea{0.45in}
\fieldline{Residual risk:}\workarea{0.45in}
\fieldline{Verification requirements:}\workarea{0.45in}
\fieldline{Related findings:}\workarea{0.35in}
\end{worksheetbox}

\subsection{Risk Characterization}
Use the organization's approved model; keep technical evidence separate from the selected scoring system.
\par\smallskip\noindent\begin{tabularx}{\linewidth}{@{}L{0.22\linewidth}L{0.17\linewidth}L{0.18\linewidth}Y@{}}
\toprule
Model / scheme & Severity & Likelihood / exploitability & Business impact / notes \\
\midrule
CVSS (if used) & & & \\
Organizational risk matrix & & & \\
Other & & & \\
\bottomrule
\end{tabularx}
\clearpage

\section{Remediation Support}
\begin{longtable}{@{}L{0.27\linewidth}L{0.67\linewidth}@{}}
\toprule
Topic & Record \\
\midrule
Developer questions & \\
Proposed fix & \\
Remediation review & \\
Regression concerns & \\
Retest result & \\
Operational mitigation & \\
Accepted risk / owner & \\
Unresolved issues & \\
Final disposition & \\
\bottomrule
\end{longtable}
\clearpage

\section{Assessment Closeout}
Avoid a simplistic overall grade. Summarize evidence, recurring patterns, limits, and residual uncertainty.
\begin{worksheetbox}{Closeout Summary}
\fieldline{Coverage summary:}\workarea{0.7in}
\fieldline{Significant security themes:}\workarea{0.7in}
\fieldline{Recurring implementation patterns:}\workarea{0.6in}
\fieldline{Architecture / design observations:}\workarea{0.6in}
\fieldline{Unresolved hypotheses:}\workarea{0.55in}
\fieldline{Known limitations:}\workarea{0.55in}
\fieldline{Areas requiring deeper analysis:}\workarea{0.55in}
\fieldline{Recommendations for future reviews:}\workarea{0.55in}
\fieldline{Remediation verification status:}\workarea{0.4in}
\end{worksheetbox}
\clearpage

\section{One-Page Execution Checklist}
\begin{enumerate}
\item Define review intent, available time, access mode, and evidence quality.
\item Build the first application model: purpose, assets, entry points, components, security expectations, trust boundaries, dependencies.
\item Select top-down, bottom-up, or hybrid structure and record switch conditions.
\item Execute bounded PLAN--WORK--REFLECT iterations with explicit exit criteria.
\item Maintain a hypothesis register, coverage register, and reconstructable working papers.
\item Convert validated issues into evidence-backed findings with prerequisites and remediation guidance.
\item Support remediation and retest; document residual risk and assessment limitations.
\end{enumerate}
\end{document}
